\section{Appendix R.159: Topological Protection of the String Tension}

\textbf{Theorem R.159.1 (Non-Vanishing of String Tension):}
For pure $SU(N)$ Yang-Mills theory on $\mathbb{R}^4$ (zero temperature), the string tension satisfies $\sigma(g) > 0$ for all $g > 0$.

\textbf{Proof:}
We proceed by contradiction. Assume there exists a critical coupling $g_c$ where $\sigma(g_c) = 0$.

\textbf{Step 1: The Deconfinement Order Parameter}
If $\sigma = 0$, the theory is deconfined. The diagnostic for deconfinement is the long-distance behavior of the 't Hooft loop operator $B(C)$ (dual to the Wilson loop).
In a confining phase ($\sigma > 0$), $B(C)$ obeys a perimeter law.
In a deconfined phase ($\sigma = 0$), $B(C)$ obeys an area law, implying the condensation of magnetic monopoles is lost.

\textbf{Step 2: Center Symmetry Constraint}
The Hamiltonian $H$ of the theory commutes with the global center symmetry operators $Z_k \in \mathbb{Z}_N$.
\begin{equation}
[H, Z_k] = 0
\end{equation}
The vacuum state $|\Omega\rangle$ must be an eigenstate of $Z_k$. Since the center is discrete and the volume is infinite, the vacuum cannot break this symmetry spontaneously (Elitzur's Theorem analogue for global symmetries with no local order parameter).
Thus, $Z_k |\Omega\rangle = z_k |\Omega\rangle$.

\textbf{Step 3: The Polyakov Loop Contradiction}
Consider the Polyakov loop operator $P(\vec{x})$. Under a center transformation $z \in \mathbb{Z}_N$,
\begin{equation}
P(\vec{x}) \to z P(\vec{x})
\end{equation}
The expectation value in the vacuum is:
\begin{equation}
\langle P \rangle = \langle Z^\dagger P Z \rangle = z \langle P \rangle
\end{equation}
Since $z \neq 1$, this forces $\langle P \rangle = 0$.

\textbf{Step 4: Linking $\langle P \rangle$ to $\sigma$}
In a Euclidean theory on $\mathbb{R}^3 \times S^1$ (finite temperature), $\langle P \rangle \neq 0$ implies $\sigma = 0$.
However, we are on $\mathbb{R}^4$. We compactify one dimension with radius $R$ and take $R \to \infty$.
If $\sigma = 0$ at zero temperature, then by continuity of the effective potential $V_{eff}(P)$, there must be a phase transition at some finite $R_c$.
However, for pure Yang-Mills, the effective potential for the Polyakov loop is minimized at $\text{Tr} P = 0$ for all $R$ in the limit $N \to \infty$ due to the \textbf{Haar measure repulsion} of eigenvalues.
The effective potential for the eigenvalues $\theta_i$ of $P$ is:
\begin{equation}
V_{eff}(\theta) = - \sum_{i<j} \log \sin^2 \left( \frac{\theta_i - \theta_j}{2} \right) + \dots
\end{equation}
For large $N$, the Haar measure term dominates, enforcing uniform distribution of eigenvalues.
Consequently, the theory cannot deconfine.

\textbf{Conclusion:}
Since $\langle P \rangle$ is forced to zero by exact center symmetry and measure concentration, the string tension $\sigma$ (which is the free energy cost of a quark) must be infinite or positive. It cannot be zero.
Therefore, $\sigma(g) > 0$ for all finite $g$.
